
\SGARefTarget{sga3:xii:expose}


From the present Exposé onwards, in contrast to the preceding Exposés,
we shall use the known results on the structure of smooth algebraic
groups over an algebraically closed field \(k\), and above all Borel's
theory of affine algebraic groups, expounded in the 1956/57 Chevalley
Seminar, \emph{Classification of Algebraic Lie Groups}. Following the usage in the theory of algebraic groups, we shall
refer to that seminar by the siglum \emph{BIBLE}. In the following
Exposés we shall need, in particular, the results of \emph{BIBLE}~4,
5, 6, and~7 (the number following \emph{BIBLE} is the number of the
Exposé). Moreover, it appears that scheme theory brings no notable
simplification to Borel's theory as it is presented in \emph{BIBLE}.
For this reason it did not seem useful to reproduce that theory in the
present Seminar: our purpose here is to deduce from results known over
algebraically closed fields analogous results valid over an arbitrary
base prescheme. (The situation will be different for Chevalley's
structure theory of semisimple groups, which, it appears, is more
advantageously treated \emph{ab ovo} over an arbitrary base
prescheme.)

In the oral Exposés (which we followed in Nos.~1 to~4), we restricted
ourselves to affine group preschemes over \(S\), relying essentially on
the representability theorems of Exposé~XI, No.~4. In the present notes
(cf. Nos.~6 to~8), we show quickly how the affineness hypotheses can be
eliminated by a simpler method that does not use the most delicate
results of Exposé~XI. For further generalizations of the results
contained in the present Exposé, see also Exposés~XV and~XVI. (It is
clear that the contents of Exposés~XI, XII, XV, and~XVI ought to be
completely rewritten.)

\section{Maximal tori}
\SGARefTarget{sga3:xii:section:1}

\paragraph{\textbf{1.0.}}
\SGARefTarget{sga3:xii:1.0}
Let first \(G\) be an algebraic group over an algebraically closed
field \(k\). A \emph{maximal torus of \(G\)} is an algebraic subgroup
\(T\) of \(G\) which is a torus (here, since \(k\) is algebraically
closed, this means that it is isomorphic to a group of the form
\(\mathbf G_m^r\)) and which is maximal for this property. Notice that,
since \(k\) is perfect, \(G_{\mathrm{red}}\) is a subgroup of \(G\),
smooth over \(k\), and hence is essentially an algebraic group in the
sense of \emph{BIBLE}. Moreover, every reduced subgroup of \(G\) (such
as a torus) is automatically a subgroup of \(G_{\mathrm{red}}\).
Consequently the maximal tori of \(G\) are precisely the maximal tori
of \(G_{\mathrm{red}}\). When \(G\) is affine, and therefore
\(G_{\mathrm{red}}\) is affine, a fundamental theorem of Borel says
that any two maximal tori of \(G\) are conjugate by an element of
\[
  G(k)=G_{\mathrm{red}}(k)
\]
(\emph{BIBLE}~6, Theorem~4(c)); in particular, they have the same
dimension. We call the common dimension of the maximal tori of \(G\)
the \emph{reductive rank of \(G\)}. The assumption that \(G\) be affine
is in fact unnecessary. This follows from a theorem of Chevalley
asserting that every connected smooth algebraic group over \(k\) is an
extension of an abelian variety by an affine group; cf. No.~6. In
Nos.~1 to~4 we shall nevertheless most often restrict ourselves to
affine group preschemes over the base.

Let now \(G\) be a smooth algebraic group over \(k\), and let \(T\) be
a maximal torus of \(G\). The centralizer \(C(T)\) of \(T\) in \(G\)
will be called the \emph{Cartan subgroup of \(G\) associated with
\(T\)}. For us it is a group subprescheme defined by \SGARefLink{sga3:viii:corollary:6:7}{VIII~6.7}. Since
\(G\) is smooth over \(k\), however, its Cartan subgroup \(C\) is also
smooth over \(k\), by \SGARefLink{sga3:xi:corollary:5:3}{XI~5.3}. Thus \(C\) is in this case the unique
group subprescheme of \(G\), smooth over \(k\) (that is, an algebraic
subgroup of \(G\) in the sense of \emph{BIBLE}), for which \(C(k)\) is
the subgroup of \(G(k)\) centralizing \(T(k)\); in other words, it is
essentially the centralizer in the sense of \emph{BIBLE}.

By the conjugacy theorem just cited, the Cartan subgroups associated
with the various maximal tori are conjugate to one another and
therefore have the same dimension. We call their common dimension the
\emph{nilpotent rank of \(G\)}; it is equal to that of
\(G_{\mathrm{red}}\). If \(\rho_r(G)\) and \(\rho_n(G)\) denote the
reductive and nilpotent ranks of \(G\), then
\[
  \rho_r(G)\leq \rho_n(G),
\]
and the difference
\[
  \rho_u(G)=\rho_n(G)-\rho_r(G)=\dim(C/T)
\]
might be called, when \(G\) is affine, the \emph{unipotent rank of
\(G\)}. If \(G\) is smooth, affine, and connected, then \(C\) is a
connected nilpotent algebraic group (\emph{BIBLE}~6,
Theorem~6(a),(c)). Hence, by \emph{BIBLE}~6, Theorem~2, it is
isomorphic to the product
\[
  C_s\times C_u,
\]
where \(C_s=T\) is the original maximal torus (and \emph{a fortiori}
a maximal torus of \(C\)), while \(C_u\) is a smooth unipotent
subgroup, that is, a successive extension of groups isomorphic to
\(\mathbf G_a\) (\emph{BIBLE}~6, Theorem~1, Corollary~1, and
\emph{BIBLE}~7, Theorem~4). In this case one also has
\[
  \rho_u(G)=\dim C_u.
\]

\paragraph{\textbf{Remark 1.1.}}
\SGARefTarget{sga3:xii:1.1}
Besides the three notions of rank just specified for an affine
algebraic group, two others are useful. The first is the
\emph{semisimple rank} \(\rho_s(G)\), defined as the reductive rank of
the quotient \(G/R\), where \(R\) is the radical of \(G\). The second
is the \emph{infinitesimal rank} \(\rho_i(G)\), defined as the
nilpotent rank of the Lie algebra of \(G\), which will be defined and
studied in the following Exposé. We shall not use them in the present
Exposé. We record only the inequalities
\[
  \rho_s\leq \rho_r\leq \rho_n\leq \rho_i.
\]

\paragraph{\textbf{Lemma 1.2.}}
\SGARefTarget{sga3:xii:1.2}
\emph{Let \(G\) be an algebraic group over the algebraically closed
field \(k\), let \(T\) be an algebraic subgroup of \(G\), and let
\(k'\) be an algebraically closed extension of \(k\). Let \(G'\) and
\(T'\) be the groups obtained from \(G\) and \(T\) by base extension.
Then \(T\) is a maximal torus of \(G\) if and only if \(T'\) is a
maximal torus of \(G'\).}

This is an immediate consequence of the ``principle of finite
extension'' (EGA~IV, 9.1.1).

\paragraph{\textbf{Definition 1.3.}}
\SGARefTarget{sga3:xii:1.3}
Let \(S\) be a prescheme, let \(G\) be an \(S\)-group prescheme of
finite type, and let \(T\) be a group subprescheme of \(G\). One says
that \(T\) is a \emph{maximal torus of \(G\)} if:
\begin{enumerate}
\renewcommand{\labelenumi}{\alph{enumi})}
\item \(T\) is a torus (\SGARefLink{sga3:ix:definition:1:3}{IX~1.3}); and
\item for every \(s\in S\), if \(\bar s\) denotes the spectrum of an
  algebraic closure of \(\kappa(s)\), then \(T_{\bar s}\) is a
  maximal torus of \(G_{\bar s}\).
\end{enumerate}

\paragraph{\textbf{Remarks 1.4.}}
\SGARefTarget{sga3:xii:1.4}
It follows from \SGARefLink{sga3:xii:1.2}{1.2} that, when \(S\) is the spectrum of an
algebraically closed field, this recovers the usual definition, and
that property~\SGARefLink{sga3:xii:1.3}{1.3} is stable under arbitrary base change. Notice that,
by \SGARefLink{sga3:x:theorem:8-8}{X~8.8}, condition~(a) in \SGARefLink{sga3:xii:1.3}{Definition~1.3} may be replaced by:
\begin{quote}
\textit{a\(^{\prime}\)) \(T\) is of finite presentation and flat over
\(S\).}
\end{quote}

One should take care: a maximal torus in the sense of \SGARefLink{sga3:xii:1.3}{Definition~1.3}
is indeed maximal among the subtori of \(G\), as follows immediately
from \SGARefLink{sga3:ix:corollary:2:9}{IX~2.9}, but the converse cannot hold---say when the base \(S\) is
connected---unless \(G\) actually admits a maximal torus in the sense
of \SGARefLink{sga3:xii:1.3}{1.3}. This is not true in general, even when
\(S=\operatorname{Spec}(\mathbb Z)\) and \(G\) is ``semisimple''; see
also \SGARefLink{sga3:xii:1.6}{1.6}. We shall see in Exposé~XIV, however, that the converse is
true when \(S\) is artinian, or when \(S\) is a local scheme and \(G\)
is ``reductive'': in that case every torus of \(G\) is contained in a
maximal torus.

\paragraph{\textbf{Definition 1.5.}}
\SGARefTarget{sga3:xii:1.5}
Let \(G\) be an algebraic group over a field \(k\). The
\emph{reductive rank} (respectively the \emph{nilpotent rank}, the
\emph{unipotent rank}, and so on) of \(G\) is the reductive rank
(respectively, and so on) of \(G_{\bar k}\), where \(\bar k\) is an
algebraic closure of \(k\).

By \SGARefLink{sga3:xii:1.2}{1.2} and the compatibility of the formation of
\(\operatorname{Centr}_G(T)\) with base extension, the notions of rank
introduced in \SGARefLink{sga3:xii:1.5}{1.5} are invariant under extension of the base field.
For algebraically closed \(k\), they coincide with the notions
introduced at the beginning of the present number.

\paragraph{\textbf{Remark 1.6.}}
\SGARefTarget{sga3:xii:1.6}
It is not difficult to construct a smooth affine group scheme \(G\)
over the spectrum \(S\) of a discrete valuation ring whose generic
fiber is isomorphic to \(\mathbf G_m\) and whose special fiber is
isomorphic to \(\mathbf G_a\). Such a \(G\) contains no torus other
than the trivial torus \(T\), reduced to the identity subgroup, and
\(T\) is plainly not a maximal torus. More precisely, in the special
fiber \(G_0=\mathbf G_{a,k}\), the torus \(T_0\) is indeed maximal,
whereas in the generic fiber \(G_1=\mathbf G_{m,K}\), the torus
\(T_1\) is no longer maximal; here \(k\) is the residue field and
\(K\) the fraction field of the valuation ring. This example also
shows that the reductive rank of \(G_s\), for \(s\in S\), is not a
continuous function of \(s\). Nevertheless, one has the following
results.

\paragraph{\textbf{Theorem 1.7.}}
\SGARefTarget{sga3:xii:1.7}
\emph{Let \(S\) be a prescheme and let \(G\) be a smooth affine
\(S\)-group prescheme. For every \(s\in S\), set
\[
  \rho_r(s)=\rho_r(G_s),\qquad
  \rho_n(s)=\rho_n(G_s),
\]
the reductive and nilpotent ranks of \(G_s\) as in \SGARefLink{sga3:xii:1.5}{1.5}. Then:}

\begin{enumerate}
\renewcommand{\labelenumi}{\alph{enumi})}
\item\SGAInlineRefTarget{sga3:xii:1.7:item:a} \emph{The function \(\rho_r\) on \(S\) is lower
  semicontinuous, while \(\rho_n\) is upper semicontinuous. Hence
  \(\rho_u=\rho_n-\rho_r\) is upper semicontinuous.}

\item\SGAInlineRefTarget{sga3:xii:1.7:item:b} \emph{The following conditions, each stable under arbitrary
  base change, are equivalent:}
  \begin{enumerate}
  \renewcommand{\labelenumii}{(\roman{enumii})}
  \item \emph{the function \(\rho_r\) on \(S\) is locally constant;}
  \item \emph{\(G\) admits a maximal torus locally for the étale
    topology;}
  \item[\textup{(ii bis)}] \emph{\(G\) admits a maximal torus locally
    for the faithfully flat quasi-compact topology.}
  \end{enumerate}

\item \emph{Let \(T_1,T_2\) be two maximal tori of \(G\), which in
  particular places us under the conditions of (b). Then \(T_1\) and
  \(T_2\) are locally conjugate for the étale topology: there exists a
  surjective étale morphism \(S'\to S\) such that the subgroups
  \((T_1)_{S'}\) and \((T_2)_{S'}\) of \(G_{S'}\) are conjugate by a
  section of \(G_{S'}\) over \(S'\).}

\item \emph{Under the conditions of (b), that is, when \(\rho_r\) is
  locally constant, \(\rho_n\) is locally constant as well, and hence
  so is \(\rho_u=\rho_n-\rho_r\).}
\end{enumerate}

\emph{Proof.}
a) Observe that for every morphism \(S'\to S\), if
\(G'=G\times_S S'\), then the functions \(\rho'_r\), and so on, on
\(S'\), defined in terms of \(G'\) as \(\rho_r\), and so on, are
defined in terms of \(G\), are obtained simply by composing the latter
functions with \(S'\to S\). When \(S'\to S\) is faithfully flat and
quasi-compact, it follows that \(\rho'\) is continuous, respectively
upper semicontinuous, respectively lower semicontinuous, if and only
if \(\rho\) is, since the Zariski topology on \(S\) is the quotient of
that on \(S'\) (SGA~1, \SGARefLink{sga3:viii:corollary:4:3}{VIII~4.3}). Consequently, the assertions in (a)
are local for the faithfully flat quasi-compact topology.

Let \(s\in S\). We wish to show that the set \(U\) of \(t\in S\) such
that
\[
  \rho_r(t)\geq \rho_r(s)
  \qquad
  \bigl(\text{respectively }\rho_n(t)\leq \rho_n(s)\bigr)
\]
is a neighborhood of \(s\). By the principle of finite extension,
there exists a finite extension \(k\) of \(\kappa(s)\) such that
\(G_k\) admits a maximal torus. There then exist an open neighborhood
\(U\) of \(s\) and a finite, surjective, locally free morphism
\(S'\to U\) such that the fiber \(S'_s\) is
\(\kappa(s)\)-isomorphic to \(\operatorname{Spec}(k)\) (cf. EGA~III,
10.3.2, where the noetherian hypothesis is plainly unnecessary).
Since \(S'\to S\) is an open morphism (SGA~1, \SGARefLink{sga3:iv:subsection:6:6}{IV~6.6}), we are reduced
to the case \(S'=S\), that is, to the case in which \(G_s\) admits a
maximal torus \(T_s\). Moreover, by \SGARefLink{sga3:xi:theorem:5:8:item:a}{XI~5.8(a)}, after possibly replacing
\(S\) once more by an étale \(S'\) over \(S\), equipped with a point
\(s'\) above \(s\), we may suppose that \(T_s\) is the fiber at \(s\)
of a subtorus \(T\) of \(G\). Then, for every \(t\in S\),
\[
  \rho_r(t)=\rho_r(G_t)
  \geq \dim T_t
  =\dim T_s
  =\rho_r(G_s)
  =\rho_r(s),
\]
which proves that \(\rho_r\) is lower semicontinuous.

On the other hand, by \SGARefLink{sga3:xi:corollary:5:3}{XI~5.3}, the functor
\[
  \operatorname{Centr}_G(T)
\]
is represented by
\[
  C=\operatorname{Centr}_G(T),\footnote{Editor's note: the text has
  been modified here to introduce the prescheme
  \(C=\operatorname{Centr}_G(T)\) representing the functor
  \(\operatorname{Centr}_G(T)\).}
\]
a closed group subprescheme of \(G\), smooth over \(S\). There is
therefore a neighborhood \(U\) of \(s\) such that \(t\in U\) implies
\[
  \dim C_t=\dim C_s=\rho_n(s).
\]
The upper semicontinuity of \(\rho_n\) now follows from the relation
\[
  \rho_n(t)\leq \dim C_t
  \qquad\text{for every }t\in S,
\]
which is itself contained in the following purely geometric lemma.

\paragraph{\textbf{Lemma 1.8.}}
\SGARefTarget{sga3:xii:1.8}
\emph{Let \(k\) be a field, let \(G\) be a smooth affine algebraic
group over \(k\), let \(T\) be a torus in \(G\), and let \(C\) be its
centralizer. Then}
\[
  \rho_n(G)\leq \dim C.
\]

Indeed, we may suppose \(k\) algebraically closed, so that \(T\) is
contained in a maximal torus \(T'\). If \(C'\) is the centralizer of
\(T'\), then \(C'\subset C\), and hence
\(\dim C'\leq \dim C\). This proves the lemma.

b) If \(\rho_r\) is locally constant, then for every torus \(T\) in
\(G\) and every \(s\in S\), if \(T_s\) is a maximal torus in \(G_s\),
there exists an open neighborhood \(U\) of \(s\) such that
\(T|_U\) is a maximal torus in \(G|_U\). Applying the argument of (a),
we see that (i) implies (ii bis). Conversely, (ii bis) implies (i),
because when \(G\) admits a maximal torus \(T\), it is clear that
\[
  \rho_r(s)=\dim T_s
\]
is a locally constant function of \(s\); and we observed at the
beginning of the proof of (a) that continuity of \(\rho_r\) is local
for the faithfully flat topology. Thus (i) is equivalent to
(ii bis), while (ii) plainly implies (ii bis); it remains to prove
the converse implication (i) \(\Rightarrow\) (ii).

For this purpose, introduce the functor \(F\) of \SGARefLink{sga3:xi:theorem:4:1}{XI~4.1},
which is a smooth separated prescheme over \(S\), and consider the
subfunctor \(\mathcal T\) of \(F\) whose value on an \(S'\) over \(S\)
is the set of maximal tori in \(G_{S'}\). By the preceding observation
that, under (i), a torus in \(G_{S'}\) which is maximal in the fiber
over a point \(s\in S'\) is maximal over an open neighborhood of
\(s\), the functor \(\mathcal T\) is represented by an open
subprescheme of \(F\), and is therefore smooth and separated over
\(S\). Since the structural morphism \(\mathcal T\to S\) is plainly
surjective, it admits a section locally for the étale topology by
\SGARefLink{sga3:xi:proposition:1:10}{XI~1.10}. This proves (i) \(\Rightarrow\) (ii).

c) This is an immediate consequence of \SGARefLink{sga3:xi:corollary:5:4:bis}{XI~5.4 bis}, taking account of
Borel's conjugacy theorem recalled at the beginning of this number.

d) By the observation at the beginning of the proof of (a), and in
view of (b), we may suppose that \(G\) admits a maximal torus \(T\).
If \(C\) is its centralizer, then \(C\) is representable and smooth
over \(S\) by \SGARefLink{sga3:xi:corollary:5:3}{XI~5.3}, so the function
\[
  s\longmapsto \rho_n(s)=\dim C_s
\]
is indeed locally constant.

This completes the proof of \SGARefLink{sga3:xii:1.7}{1.7}. We shall refer to the conditions in
\SGARefLink{sga3:xii:1.7:item:b}{1.7(b)} by saying that, in this case, \(G\) is \emph{of locally
constant reductive rank}. We record the following consequence.

\paragraph{\textbf{Corollary 1.9.}}
\SGARefTarget{sga3:xii:1.9}
\emph{Let \(G\) be as in \SGARefLink{sga3:xii:1.7}{1.7}, and let \(s\in S\) be such that
\(\rho_u(s)=0\), that is,
\(\rho_r(s)=\rho_n(s)\) (equivalently, the Cartan subgroups of
\(G_{\bar k}\) are tori, where \(\bar k\) is an algebraic closure of
\(\kappa(s)\)). Then there exists an open neighborhood \(U\) of \(s\)
on which \(\rho_r\) and \(\rho_n\) are constant. In particular, for
every \(t\in U\), the unipotent rank \(\rho_u(t)\) of \(G_t\) is
zero.}

Indeed, this follows immediately from \SGARefLink{sga3:xii:1.7:item:a}{1.7(a)} and the inequality
\(\rho_r(t)\leq \rho_n(t)\) for every \(t\in S\).

Notice also that, in the course of proving (b), we proved the
following result.

\paragraph{\textbf{Corollary 1.10.}}
\SGARefTarget{sga3:xii:1.10}
\emph{Let \(G\) be as in \SGARefLink{sga3:xii:1.7}{1.7}, and suppose that \(G\) has locally
constant reductive rank. Consider the functor
\[
  \mathcal T:(\mathbf{Sch})^\circ\longrightarrow(\mathbf{Ens})
\]
such that, for every \(S'\) over \(S\),
\[
  \mathcal T(S')=
  \{\text{maximal tori of }G_{S'}\}.
\]
Then \(\mathcal T\) is represented by a smooth, separated prescheme of
finite type over \(S\).}

It remains only to verify that \(\mathcal T\) is of finite type over
\(S\). Since the question is local for the faithfully flat
quasi-compact topology, we may suppose that \(G\) admits a maximal
torus \(T\). By \SGARefLink{sga3:xi:corollary:5:3:bis}{XI~5.3 bis} and the bis version of \SGARefLink{sga3:xi:remark:5:5}{XI~5.5},
\(\operatorname{Norm}_G(T)\) and
\(G/\operatorname{Norm}_G(T)\) are represented by preschemes
\(\operatorname{Norm}_G(T)\) and
\(G/\operatorname{Norm}_G(T)\), and \(\mathcal T\) is isomorphic to
\[
  G/\operatorname{Norm}_G(T).\footnote{Editor's note: the text has
  been modified here to introduce the preschemes
  \(\operatorname{Norm}_G(T)\) and
  \(G/\operatorname{Norm}_G(T)\).}
\]
The morphism \(G\to\mathcal T\) defined by
\(g\mapsto\operatorname{int}(g)(T)\) is surjective. Since \(G\) is
quasi-compact over \(S\), so is \(\mathcal T\), which completes the
proof.

\paragraph{\textbf{Remarks 1.11.}}
\SGARefTarget{sga3:xii:1.11}
a) The prescheme \(\mathcal T\) of \SGARefLink{sga3:xii:1.10}{1.10} will properly be called the
\emph{prescheme of maximal tori of \(G\)}. We shall see in No.~5 that
it is in fact affine over \(S\). This can be checked directly when
\(G\) is isomorphic to a closed group subprescheme of a prescheme of
the form \(\operatorname{GL}(n)\), using \SGARefLink{sga3:xi:remark:4:6}{XI~4.6} and observing that, by
the argument in the proof of \SGARefLink{sga3:xii:1.7:item:b}{1.7(b)}, \(\mathcal T\) identifies with a
subprescheme both open and closed in the prescheme \(F\) representing
the multiplicative-type subgroups of \(G\).

b) It is possible to prove \SGARefLink{sga3:xii:1.7}{1.7}, and hence \SGARefLink{sga3:xii:1.9}{1.9}, without using the
delicate results of Exposé~XI, using only the elementary results
\SGARefLink{sga3:xi:proposition:3:12}{XI~3.12} and \SGARefLink{sga3:xi:theorem:6:2}{XI~6.2} and working solely with multiplicative-type groups
finite over the base (morally, the groups \({}_nT\), where \(T\) is a
maximal torus); compare No.~7. The same observation applies to the
proof of \SGARefLink{sga3:xii:1.10}{1.10}.

We end this number with examples in which there is a unique maximal
torus.

\paragraph{\textbf{Proposition 1.12.}}
\SGARefTarget{sga3:xii:1.12}
\emph{Let \(G\) be an \(S\)-group prescheme of multiplicative type
and finite type. Then \(G\) admits a unique maximal torus, and every
torus in \(G\) is contained in this maximal torus.}

Uniqueness plainly follows from the latter assertion, which
characterizes the maximal torus as the largest subtorus of \(G\).
Uniqueness also makes existence a question local for the faithfully
flat quasi-compact topology. We may therefore suppose that \(G\) is
diagonalizable, say \(G=D_S(M)\), where \(M\) is a finitely generated
commutative group. Let \(M_0\) be the quotient of \(M\) by its torsion
subgroup. We claim that the torus
\[
  T=D_S(M_0)
\]
in \(G\) is a maximal torus and the largest subtorus. Indeed, a
subtorus \(T'\) of \(G\) is locally diagonalizable for the fpqc
topology. Thus, in order to prove \(T'\subset T\), we may suppose that
\(T'\) is diagonalizable, and hence of the form \(D_S(N)\), where
\(N\) is a free quotient of \(M\) (\SGARefLink{sga3:viii:corollary:1:4}{VIII~1.4} and \SGARefLink{sga3:viii:corollary:3:2:item:b:01}{3.2(b)}). It follows
that \(N\) is a quotient of \(M_0\), and therefore that
\(T'\subset T\). Since the construction of \(T=D_S(M_0)\) is
compatible with arbitrary base extension, this proves at the same
time that \(T\) is a maximal torus of \(G\), and completes the proof.
When \(G\) is smooth over \(S\), \SGARefLink{sga3:xii:1.12}{Proposition~1.12} has the following
generalization.

\paragraph{\textbf{Proposition 1.13.}}
\SGARefTarget{sga3:xii:1.13}
\emph{Let \(G\) be an \(S\)-group prescheme of finite presentation
over \(S\). Suppose that, locally for the faithfully flat
quasi-compact topology, \(G\) admits a central maximal torus. Then
\(G\) admits globally a unique maximal torus, and it is the largest subtorus of \(G\).}

Here again, uniqueness follows trivially from the last assertion and
makes existence a local question for the fpqc topology. We may
therefore suppose that \(G\) admits a central maximal torus \(T\). We
shall show that every torus \(R\) of \(G\) is contained in \(T\). This
follows from the next lemma.

\paragraph{\textbf{Lemma 1.14.}}
\SGARefTarget{sga3:xii:1.14}
\emph{Let \(G\) be an \(S\)-group prescheme of finite presentation
over \(S\), let \(T\) be a maximal torus of \(G\), let \(R\) be a
subtorus of \(G\), and suppose that \(T\) and \(R\) commute. Then
\(R\subset T\).}

Indeed, since \(R\) and \(T\) commute, the morphism
\[
  R\times_S T\longrightarrow G,\qquad (r,t)\longmapsto r\cdot t,
\]
is a group homomorphism. By \SGARefLink{sga3:ix:theorem:6:8}{IX~6.8} it therefore has an image subgroup
\(T'\) in \(G\), which is a multiplicative-type group quotient of
\(R\times_S T\), hence a torus, and plainly contains \(T\). Since
\(T\) is a maximal torus, \(T'=T\) by \SGARefLink{sga3:xii:1.4}{1.4}, and therefore
\(R\subset T\). This proves \SGARefLink{sga3:xii:1.14}{1.14} and hence \SGARefLink{sga3:xii:1.13}{1.13}. In particular, using
\SGARefLink{sga3:xii:1.7:item:b}{1.7(b)}, we obtain:

\paragraph{\textbf{Corollary 1.15.}}
\SGARefTarget{sga3:xii:1.15}
\emph{Let \(G\) be a commutative, smooth, affine \(S\)-group
prescheme of locally constant reductive rank. Then \(G\) admits a
unique maximal torus, and this torus contains every subtorus of
\(G\).}

\paragraph{\textbf{Corollary 1.16.}}
\SGARefTarget{sga3:xii:1.16}
\emph{Let \(G\) be a smooth affine \(S\)-group prescheme. Suppose
that, for every \(s\in S\), if \(\bar s\) denotes the spectrum of an
algebraic closure \(\bar k\) of \(\kappa(s)\), the geometric fiber
\(G_{\bar s}\) is a connected nilpotent algebraic group (in the sense
of \emph{BIBLE}, that is, the group \(G(\bar k)\) of its
\(\bar k\)-valued points is nilpotent). Suppose moreover that \(G\)
has locally constant reductive rank. Then \(G\) admits a unique
maximal torus \(T\); moreover, \(T\) is central and is the largest
subtorus of \(G\).}

By 1.7(b), \(G\) admits a maximal torus locally for the fpqc topology,
and by \SGARefLink{sga3:xii:1.13}{1.13} we are reduced to proving that a maximal torus of \(G\) is
central. By \SGARefLink{sga3:ix:theorem:5:6:item:b}{IX~5.6(b)}, we are reduced to the case in which \(S\) is
the spectrum of a field, which we may suppose algebraically closed.
Then \(T(k)\) lies in the center of \(G(k)\) by \emph{BIBLE}~6,
\SGARefLink{sga3:xii:section:2}{Theorem~2}. This implies that \(T\) lies in the center of \(G\), because
\(\operatorname{Centr}_G(T)\) is a closed subscheme of \(G\)
(VIII~6.6) containing all the points of \(G(k)\), and is therefore
equal to \(G\) by the Nullstellensatz, since \(G\) is reduced.

Finally, for later reference, we record the following elementary
proposition, which we have already used implicitly.

\paragraph{\textbf{Proposition 1.17.}}
\SGARefTarget{sga3:xii:1.17}
\emph{Let \(G\supset H\supset T\) be group preschemes of finite type
over \(S\). The following conditions are equivalent:}
\begin{enumerate}
\renewcommand{\labelenumi}{(\roman{enumi})}
\item \emph{\(T\) is a maximal torus of \(G\);}
\item \emph{\(T\) is a maximal torus of \(H\), and for every
  \(s\in S\) one has equality of reductive ranks
  \[
    \rho_r(G_s)=\rho_r(H_s).
  \]}
\end{enumerate}

\section{The Weyl group}
\SGARefTarget{sga3:xii:section:2}

Let first \(k\) be an algebraically closed field, and let \(G\) be a
smooth affine algebraic group over \(k\). If \(T\) is a maximal torus,
\(C\) its centralizer, and \(N\) its normalizer, then by \SGARefLink{sga3:xi:proposition:5:9}{XI~5.9} these
are smooth closed subgroups of \(G\), and \(C\) is an open subgroup of
\(N\). Thus
\[
  W=N/C
\]
is a finite étale group over \(k\), and is consequently determined by
the group \(W(k)\) of its \(k\)-valued points:
\[
  W=W(k)_k
\]
(the constant group defined by the ordinary finite group \(W(k)\)).
The finite group \(W(k)\) will be called the \emph{geometric Weyl
group}, or simply the \emph{Weyl group} when no confusion is likely,
of \(G\) relative to \(T\). By Borel's conjugacy theorem, the Weyl
groups relative to the various maximal tori are isomorphic to one
another; this is why one sometimes speaks of ``the'' Weyl group of
\(G\), without specifying a maximal torus. Since the formation of
\(C\), \(N\), and \(N/C\) commutes with arbitrary base extension, if
\(k'\) is an algebraically closed extension of \(k\), the geometric
Weyl group of \(G_{k'}\) relative to \(T_{k'}\) is canonically
isomorphic to that of \(G\) relative to \(T\). Consequently, ``the''
geometric Weyl group of \(G\), properly speaking an isomorphism class
of ordinary finite groups, coincides with that of \(G_{k'}\).

This permits us, when \(G\) is a smooth affine algebraic group over an
arbitrary field \(k\), to speak of the \emph{geometric Weyl group of
\(G\)} as the isomorphism class of the geometric Weyl group of
\(G_{k'}\), where \(k'\) is any algebraically closed extension of
\(k\). When \(G\) admits a maximal torus \(T\), we can form
\[
  C=\operatorname{Centr}_G(T),\qquad
  N=\operatorname{Norm}_G(T),\qquad
  W=N/C
\]
as above. The latter is a finite étale group over \(k\), called the
\emph{Weyl group of \(G\) relative to \(T\)}. The geometric Weyl group
is then simply the class of the group of points of \(W\) with values
in any algebraically closed extension \(k'\) of \(k\). In general,
knowledge of the geometric Weyl group \(W(k')\) plainly no longer
suffices to reconstruct the algebraic group \(W\): one must also know
the action on \(W(k')\) of the Galois group of the separable algebraic
closure of \(k\) in \(k'\).

Finally, let \(G\) be a group prescheme over an arbitrary base \(S\),
smooth and affine over \(S\), and let \(T\) be a maximal torus of
\(G\). By \SGARefLink{sga3:xi:proposition:5:9}{XI~5.9} we may again form the group
\[
  W(T)=N(T)/C(T),
\]
which is an étale, separated, quasi-finite \(S\)-group prescheme. Its
geometric fibers, relative to algebraic closures of the residue
fields \(\kappa(s)\), \(s\in S\), are the geometric Weyl groups of the
fibers \(G_s\). Thus \SGARefLink{sga3:xi:corollary:5:10}{XI~5.10} gives information about the variation of
these groups with \(s\in S\). We can make this information more
precise and generalize it as follows.

\paragraph{\textbf{Theorem 2.1.}}
\SGARefTarget{sga3:xii:2.1}
\emph{Let \(S\) be a prescheme and let \(G\) be a smooth affine
\(S\)-group prescheme. For every \(s\in S\), let \(w(s)\) be the
geometric Weyl group of \(G_s\), hence an isomorphism class of finite
groups. On the set \(E\) of such classes, introduce the following
preorder: \(w\leq w'\) if and only if \(w\) and \(w'\) can be
represented by finite groups \(W\) and \(W'\) such that \(W\) is
isomorphic to a quotient of a subgroup of \(W'\). With this
convention:}
\begin{enumerate}
\renewcommand{\labelenumi}{\alph{enumi})}
\item\SGAInlineRefTarget{sga3:xii:2.1:item:a} \emph{The function \(s\mapsto w(s)\) from \(S\) to \(E\) is
  lower semicontinuous.}
\item\SGAInlineRefTarget{sga3:xii:2.1:item:b} \emph{Suppose that \(G\) has locally constant reductive rank.
  Then the following conditions are equivalent:}
  \begin{enumerate}
  \renewcommand{\labelenumii}{(\roman{enumii})}
  \item \emph{the function \(s\mapsto w(s)\) is locally constant;}
  \item[\textup{(i bis)}] \emph{the function
    \(s\mapsto\operatorname{card}w(s)\) is locally constant;}
  \item \emph{locally for the étale topology (or merely for the fpqc
    topology), there exists a maximal torus \(T\) such that \(W(T)\)
    is finite over its base;}
  \item[\textup{(ii bis)}] \emph{for every \(S'\) over \(S\) and
    every maximal torus \(T\) of \(G_{S'}\), the associated Weyl group
    \(W(T)\) is finite over \(S'\).}
  \end{enumerate}
\end{enumerate}

\emph{Proof.}
a) Arguing as in \SGARefLink{sga3:xii:1.7:item:a}{1.7(a)}, in order to prove that for every \(s\in S\)
there exists an open neighborhood \(U\) of \(s\) such that
\(t\in U\) implies \(w(t)\geq w(s)\), we are reduced to the case in
which there is a torus \(R\) in \(G\) such that \(R_s\) is a maximal
torus in \(G_s\). Put
\[
  W(R)=\operatorname{Norm}_G(R)/\operatorname{Centr}_G(R).
\]
As in \SGARefLink{sga3:xi:proposition:5:9}{XI~5.9}, this is an étale, separated, quasi-finite group
prescheme over \(S\). For \(t\in S\), let \(w'(t)\in E\) be its
geometric fiber at \(t\). Since \(R_s\) is a maximal torus in \(G_s\),
and the formation of \(\operatorname{Norm}\),
\(\operatorname{Centr}\), and
\(\operatorname{Norm}/\operatorname{Centr}\) commutes with arbitrary
base extension, in particular with passage to fibers, one has
\[
  w(s)=w'(s).
\]
We claim moreover that for \(t\) near \(s\),
\[
  w(t)\geq w'(t)
  \qquad\text{and}\qquad
  w'(t)\geq w'(s),
\]
which will suffice to prove (a). These two inequalities are contained
in the following two lemmas.

\paragraph{\textbf{Lemma 2.2.}}
\SGARefTarget{sga3:xii:2.2}
\emph{Let \(S\) be a prescheme, and let \(W\) be an étale,
separated, quasi-finite \(S\)-group prescheme. For every \(s\in S\),
let \(f(s)\) be the class of the geometric fiber of \(W\) at \(s\),
an element of the ordered set \(E\) of isomorphism classes of finite
groups introduced in \SGARefLink{sga3:xii:2.1}{2.1}. Then the function
\(f:S\to E\) is lower semicontinuous. It is constant in a
neighborhood of \(s\in S\) if and only if the function
\(s\mapsto\operatorname{card}f(s)\) is constant there, and this holds
if and only if \(W\) is finite over \(S\) above some open
neighborhood \(U\) of \(s\).}

This result refines, in group-theoretic terms, the result used in the
proof of \SGARefLink{sga3:xi:corollary:5:10}{XI~5.10}. We give only a sketch of the proof, which is entirely
standard. As usual, we are reduced to the case in which \(S\) is
affine and noetherian. The function \(f\) is immediately seen to be
constructible (EGA \(0_{\mathrm{III}}\), 9.3.1 and 9.3.2, together
with the sorites of EGA~IV, 9), and EGA
\(0_{\mathrm{III}}\), 9.3.4 reduces semicontinuity to proving that if
\(t\) is a generalization of \(s\), then
\[
  f(t)\geq f(s).
\]
By EGA~II, 7.1.7, this reduces us to the case in which \(S\) is the
spectrum of a discrete valuation ring, which may be supposed complete
with algebraically closed residue field. Let \(s\) be the closed point
of \(S\). Since \(W\) is étale and separated over \(S\), it contains
a subscheme \(W'\), both open and closed and finite over \(S\), such
that \(W'_s=W_s\) (EGA~II, 6.2.6); it is immediate here that \(W'\)
is a subgroup of \(W\). Moreover, since \(W'\) is finite étale over
\(S=\operatorname{Spec}(V)\), with \(V\) complete and with
algebraically closed residue field, it is a constant group, hence of
the form \(A_S\), where
\[
  A=W'(\kappa(s))=W(\kappa(s))
\]
has class \(f(s)\). If \(B\) is the geometric fiber of \(W\) at the
generic point \(t\) of \(S\), there is therefore a canonical
monomorphism \(A\to B\), proving \(f(s)\leq f(t)\). Notice that this
argument actually proves semicontinuity for an order on \(E\) finer
than the one specified in \SGARefLink{sga3:xii:2.1}{2.1}.

The equivalence between constancy of \(f\) at \(s\) and constancy of
\(s\mapsto\operatorname{card}f(s)\) follows from the fact that, for
\(w,w'\in E\),
\[
  w\leq w'
  \quad\text{and}\quad
  \operatorname{card}w=\operatorname{card}w'
  \quad\Longrightarrow\quad
  w=w'.
\]
The equivalence of this condition with finiteness of \(W\) over a
neighborhood of \(s\) is independent of the group structure on
\(W\); it was noted after \SGARefLink{sga3:xi:corollary:5:10}{XI~5.10}, and also follows readily from the
preceding argument using the valuative criterion of properness
(EGA~II, 7.3.8).

\paragraph{\textbf{Lemma 2.3.}}
\SGARefTarget{sga3:xii:2.3}
\emph{Let \(G\) be a smooth affine algebraic group over an
algebraically closed field \(k\), and let \(R\subset T\) be two
subtori. Let \(W(R)\) and \(W(T)\) be the two associated finite groups
of \SGARefLink{sga3:xi:proposition:5:9}{XI~5.9}, obtained as the quotient of the normalizer by the
centralizer. Then \(W(R)\) is isomorphic to a quotient of a subgroup
of \(W(T)\).}

Indeed, consider the diagram
\[
\begin{tikzcd}[column sep=large,row sep=large]
T \arrow[r, hook] &
C(T) \arrow[r, hook] \arrow[d, hook] &
C(R) \arrow[d, hook] \\
& N(T)\cap N(R) \arrow[r, hook] \arrow[d, hook] &
N(R) \\
& N(T) &
\end{tikzcd}
\]
Then \((N(T)\cap N(R))/C(T)\) is a subgroup of
\[
  W(T)=N(T)/C(T),
\]
and there is an evident homomorphism
\[
  (N(T)\cap N(R))/C(T)
  \longrightarrow
  W(R)=N(R)/C(R).
\]
Everything reduces to proving that this last homomorphism is
surjective, that is, that for every \(k\)-valued point \(g\) of
\(N(R)\), there exists a \(k\)-valued point \(c\) of \(C(R)\) such
that \(cg\) normalizes \(T\), or equivalently such that
\[
  \operatorname{int}(c)\bigl(\operatorname{int}(g)T\bigr)=T.
\]
For this, it suffices to observe that
\(\operatorname{int}(g)T\) is a torus of \(N(R)\), hence of \(C(R)\),
which is an open subgroup of \(N(R)\). The tori \(T\) and
\(\operatorname{int}(g)T\) are maximal tori of \(C(R)\), since they
are maximal in \(G\), and the assertion follows from Borel's conjugacy
theorem.

This proves \SGARefLink{sga3:xii:2.3}{2.3}, and hence \SGARefLink{sga3:xii:2.1:item:a}{2.1(a)}.

b) We have already observed that (i) and (i bis) are plainly
equivalent. They imply (ii bis), by the converse part of \SGARefLink{sga3:xii:2.2}{2.2} or,
equivalently, \SGARefLink{sga3:xi:corollary:5:10}{XI~5.10}. Condition (ii bis) implies (ii) by \SGARefLink{sga3:xii:1.7:item:b}{1.7(b)}.
Finally, (ii) implies (i bis): as in \SGARefLink{sga3:xii:1.7:item:a}{1.7(a)}, condition (i) is local
for the fpqc topology, so we may suppose that \(G\) admits a maximal
torus \(T\) such that \(W(T)\) is finite over \(S\), and \SGARefLink{sga3:xi:corollary:5:10}{XI~5.10}
again gives the conclusion. This completes the proof of \SGARefLink{sga3:xii:2.1}{2.1}.

\section{Cartan subgroups}
\SGARefTarget{sga3:xii:section:3}

\paragraph{\textbf{Definition 3.1.}}
\SGARefTarget{sga3:xii:3.1}
Let \(G\) be a smooth group prescheme of finite type over a prescheme
\(S\). A \emph{Cartan subgroup of \(G\)} is a group subprescheme
\(C\) of \(G\), smooth over \(S\), such that for every \(s\in S\), if
\(\bar s\) denotes the spectrum of an algebraic closure of
\(\kappa(s)\), then \(C_{\bar s}\) is a Cartan subgroup, in the sense
of No.~1, of \(G_{\bar s}\).

It is immediate that if \(C\) is a Cartan subgroup of \(G\), then for
every \(S'\) over \(S\), \(C_{S'}\) is a Cartan subgroup of
\(G_{S'}\). If \(S\) is the spectrum of an algebraically closed field,
this recovers the notion recalled in No.~1. Finally, the property that
a group subprescheme \(C\) of \(G\) be a Cartan subgroup is plainly
local for the faithfully flat quasi-compact topology.

\paragraph{\textbf{Theorem 3.2.}}
\SGARefTarget{sga3:xii:3.2}
\emph{Let \(G\) over \(S\) be as in \SGARefLink{sga3:xii:3.1}{3.1}, and suppose that \(G\) is
affine over \(S\) and has locally constant reductive rank. Then the
map
\[
  T\longmapsto\operatorname{Centr}_G(T)
\]
induces a bijection from the set of maximal tori of \(G\) onto the set
of Cartan subgroups\footnote{The proof given here in fact proves
\SGARefLink{sga3:xii:3.2}{Theorem~3.2} only for closed Cartan subgroups of \(G\). However, \SGARefLink{sga3:xii:7.1:item:a}{7.1(a)}
establishes \SGARefLink{sga3:xii:3.2}{Theorem~3.2} as stated, and implies that the Cartan
subgroups of \(G\) are closed.} of \(G\). If \(C\) corresponds to
\(T\), then \(T\) is the unique maximal torus of \(C\).}

If \(T\) is a maximal torus of \(G\), the functor
\(\operatorname{Centr}_G(T)\) is represented, by \SGARefLink{sga3:xi:corollary:5:3}{XI~5.3} or \SGARefLink{sga3:xii:6.2}{6.2}, by
a closed smooth subprescheme
\[
  C=\operatorname{Centr}_G(T)
\]
of \(G\), and the definitions show that \(C\) is a Cartan subgroup
of \(G\). Moreover, \(T\) is plainly a maximal torus of \(C\), and,
being central, it is the unique maximal torus of \(C\) by \SGARefLink{sga3:xii:1.13}{1.13}.
Thus the map \(T\mapsto\operatorname{Centr}_G(T)\) from the set of
maximal tori to the set of Cartan subgroups is injective; it remains
to prove surjectivity.

Let \(C\) be a Cartan subgroup of \(G\). We shall prove that
\(C=\operatorname{Centr}_G(T)\) for a maximal torus \(T\) of \(G\).
It is enough to find a maximal torus \(T\) of \(C\), for then \(T\)
is a maximal torus of \(G\): for every \(s\in S\), the groups
\(G_s\) and \(C_s\) have the same reductive rank. By \SGARefLink{sga3:ix:theorem:5:6:item:b}{IX~5.6(b)},
\(T\) lies in the center of \(C\), so
\[
  C\subset C'=\operatorname{Centr}_G(T).
\]
The smooth subgroup \(C\) of the smooth group \(C'\) over \(S\)
coincides with \(C'\) fiber by fiber; hence \(C=C'\).

Since \(G\) has locally constant reductive rank, so does \(C\).
Therefore \(C\) admits a maximal torus locally for the étale topology
by \SGARefLink{sga3:xii:1.7:item:b}{1.7(b)}; that torus is central by the preceding argument, and \SGARefLink{sga3:xii:1.13}{1.13}
then shows that \(C\) itself admits a maximal torus. This completes
the proof.

\paragraph{\textbf{Corollary 3.3.}}
\SGARefTarget{sga3:xii:3.3}
\emph{Let \(G\) be a smooth affine group prescheme over \(S\), of
locally constant reductive rank, and let
\[
  \mathcal C:(\mathrm{Sch})^\circ_{/S}\longrightarrow(\mathrm{Ens})
\]
be the functor defined by
\[
  \mathcal C(S')=
  \{\text{Cartan subgroups of }G_{S'}\}.
\]
Then \(\mathcal C\) is isomorphic to the functor \(\mathcal T\) of
\SGARefLink{sga3:xii:1.10}{1.10}, and is therefore represented by the same smooth separated
prescheme of finite type over \(S\).}

\paragraph{\textbf{Corollary 3.4.}}
\SGARefTarget{sga3:xii:3.4}
\emph{Under the hypotheses of \SGARefLink{sga3:xii:3.2}{3.2}, if
\(C=\operatorname{Centr}_G(T)\), then
\[
  \operatorname{Norm}_G(C)=\operatorname{Norm}_G(T).
\]}

\paragraph{\textbf{Remark 3.5.}}
\SGARefTarget{sga3:xii:3.5}
When \(G\) does not have locally constant reductive rank, it may
nevertheless admit Cartan subgroups. For example, if \(G\) has
connected nilpotent fibers, then \(G\) is a Cartan subgroup of itself,
although it need not have locally constant reductive rank; see \SGARefLink{sga3:xii:1.6}{1.6}.
In Exposé XV we develop the theory of Cartan subgroups without
assuming either that \(G\) is affine over \(S\) or that its reductive
rank is locally constant, using the theory of regular elements from
the next Exposé. See also Nos.~6 and 7 below for the removal of the
affineness hypothesis.
